% ============================================================
% RESULT PRESENTATION TEMPLATE
% For EDC Book 2 - Use for ALL major results
% Date: 2026-01-30
% ============================================================

% ============================================================
% TEMPLATE USAGE:
% Copy this entire block for each major result
% Replace [PLACEHOLDERS] with actual content
% Delete instruction comments before final compilation
% ============================================================

\subsection{[Result Name]}
\label{sec:[result_label]}

% ═══════════════════════════════════════════════════════════
% SECTION 1: GEOMETRIC DERIVATION
% ═══════════════════════════════════════════════════════════

\subsubsection{Derivation Within Model}

\textbf{Geometric setup} [P]:

[Describe 5D structure, boundary conditions, symmetries, etc.]

Example:
% From 5D membrane with Z_6 crystallographic symmetry,
% Y-junction configuration with three strings meeting at angles θ_i.
% Boundary conditions: frozen at scale δ = ℏ/(2m_p c).

\textbf{Mathematical calculation} yields [\textcolor{blue}{Dc:Sym}]:
\begin{equation}
  [\text{quantity}] = [\text{exact formula with } \pi, \text{ etc.}]
  \label{eq:[result_formula]}
\end{equation}

Example:
% \frac{m_p}{m_e} = 6\pi^5

\textbf{Numerical evaluation}:
\begin{equation}
  [\text{quantity}] = [\text{decimal value to many digits}]
\end{equation}

Example:
% \frac{m_p}{m_e} = 1836.1181148335...

\textbf{Epistemic status}: [\textcolor{blue}{Dc:Sym}] — Conditionally derived (symbolic exact)

Conditional on:
\begin{enumerate}
\item 5D postulates [P]: [list specific postulates]
  % Example: Y-junction topology, Z_6 symmetry
\item Reduction dictionary: "[\text{This 5D quantity}] $\equiv$ [\text{PDG observable}]"
  % Example: "E_5D/c² ≡ mass_3D"
\item Normalization: $C_{\mathrm{red}}$ (currently unconstrained [Open])
\item Approximations: [list, e.g., "tree level", "leading order in $\alpha$"]
\end{enumerate}

% ═══════════════════════════════════════════════════════════
% SECTION 2: EMPIRICAL COMPARISON
% ═══════════════════════════════════════════════════════════

\subsubsection{Comparison to Measurement}

\textbf{Baseline data} [BL]: Table~\ref{tab:baseline_constants}

PDG/CODATA value:
\begin{equation}
  [\text{quantity}]_{\text{exp}} = [\text{value}] \pm [\text{uncertainty}]
\end{equation}

Example:
% \left(\frac{m_p}{m_e}\right)_{\mathrm{exp}} = 1836.15267343 \pm 1.1 \times 10^{-8}

\textbf{Agreement calculation}:
\begin{align}
\text{EDC prediction:} \quad &[\text{value}_{\text{pred}}] \\
\text{Measurement:} \quad &[\text{value}_{\text{exp}}] \pm [\text{error}] \\
\text{Absolute difference:} \quad &|\Delta| = [\text{number}] \\
\text{Relative difference:} \quad &[\text{percentage}\%] = [\text{ppm}]\,\text{ppm}
\end{align}

Example:
% Absolute: |1836.118 - 1836.153| = 0.035
% Relative: 0.035/1836.153 = 0.0019% = 19 ppm

% ═══════════════════════════════════════════════════════════
% SECTION 3: ERROR BUDGET
% ═══════════════════════════════════════════════════════════

\subsubsection{Error Budget}

\begin{table}[h]
\centering
\begin{tabular}{llcc}
\hline
\textbf{Correction Source} & \textbf{Mechanism} & \textbf{Estimate} & \textbf{Status} \\
\hline
EM corrections & $O(\alpha)$ loops & $\sim X\%$ & [Open] \\
RG running & $\beta$-functions & $\sim Y\%$ & [Open] \\
Finite-size & $(r/R)^n$ & $\sim Z\%$ & Negligible/[Open] \\
$C_{\mathrm{red}}$ & Normalization & Unknown & \textbf{[Open]} \\
[Other] & [Mechanism] & $\sim W\%$ & [Open] \\
\hline
\textbf{Total expected} & & \textbf{$A{-}B\%$} & — \\
\textbf{Observed} & & \textbf{$C\%$} & [BL] \\
\hline
\end{tabular}
\caption{Error budget for [result name]}
\label{tab:error_budget_[result]}
\end{table}

Example:
% EM corrections: O(α) ~ 0.1% [Open]
% Total expected: 0.1-1%
% Observed: 0.002%

\textbf{Interpretation}:
\begin{itemize}
\item Observed $C\%$ is within expected range $[A, B]\%$
\item Largest uncertainty: $C_{\mathrm{red}}$ normalization [Open]
\item Agreement is \textbf{strong evidence} for geometric mechanism
\item NOT conclusive proof until error budget closed
\end{itemize}

% ═══════════════════════════════════════════════════════════
% SECTION 4: SENSITIVITY ANALYSIS
% ═══════════════════════════════════════════════════════════

\subsubsection{Sensitivity Analysis}

\textbf{Explored parameter space} (Section~\ref{sec:epistemic_standard}):

\begin{itemize}
\item $[\text{param}_1]$: $\pm X\%$ variation
  $\to$ result shifts by $\pm Y\%$ (linear/quadratic/negligible)
  
\item $[\text{param}_2]$: factor of 2 variation
  $\to$ result changes by $Z\%$
  
\item $C_{\mathrm{red}}$: unconstrained [Open]
  $\to$ \textbf{dominant uncertainty}
\end{itemize}

Example:
% σ: ±10% → no change (cancels in ratio) ✓
% ℓ_p/r_e: ±1% → ±1% shift
% C_red: unknown → dominant

\textbf{Stability conclusion}:
Within explored ranges (excluding $C_{\mathrm{red}}$), prediction stable to $\pm W\%$.

\textbf{Constraint from agreement}:
Current agreement implies $C_{\mathrm{red}} \in [0.99, 1.01]$ (95\% CL, empirical).

% ═══════════════════════════════════════════════════════════
% SECTION 5: STATUS SUMMARY BOX
% ═══════════════════════════════════════════════════════════

\subsubsection{Status Summary}

\begin{tcolorbox}[colback=blue!5,colframe=blue!75!black,title=Result Status]

\begin{tabular}{ll}
\textbf{Mathematical precision:} & [Dc:Sym] (symbolic exact) \\
\textbf{Physical framework:} & Incomplete (reduction not fully known) \\
\textbf{Empirical agreement:} & $C\%$ (within expected $[A, B]\%$) \\
\textbf{Epistemic tag:} & [Dc:Sym] \\
\textbf{Confidence level:} & High, bounded by stated incompleteness \\
\textbf{Dominant uncertainty:} & $C_{\mathrm{red}}$ normalization [Open] \\
\end{tabular}

\end{tcolorbox}

% ═══════════════════════════════════════════════════════════
% SECTION 6: FALSIFICATION CRITERIA
% ═══════════════════════════════════════════════════════════

\subsubsection{Tension and Falsification}

\textbf{Current status}:
\begin{itemize}
\item NOT in tension (observed $C\%$ within expected $[A,B]\%$)
\item Error budget NOT closed ([Open] corrections remain)
\end{itemize}

\textbf{Would be in tension if}:
\begin{itemize}
\item Future measurements shifted by $> [threshold]\%$
\item After calculating [Open] corrections, residual $> 3\sigma$
\end{itemize}

\textbf{Would be falsified if}:
\begin{itemize}
\item Geometric assumption [X] proven incompatible with [experiment Y]
\item Multiple predictions simultaneously fail after error budget closure
\item No plausible $C_{\mathrm{red}}$ can restore agreement
\end{itemize}

% ═══════════════════════════════════════════════════════════
% SECTION 7: PATH FORWARD
% ═══════════════════════════════════════════════════════════

\subsubsection{Path to Higher Certainty}

To strengthen this result:
\begin{enumerate}
\item Complete brane reduction theory (derive $C_{\mathrm{red}}$ from 5D action)
\item Calculate all [Open] corrections explicitly
\item Verify stability under broader parameter variations
\item Independent derivation of dictionary mapping
\item Cross-check with complementary approaches
\end{enumerate}

Completing items (1-2) would upgrade [Dc] $\to$ [Der].

% ═══════════════════════════════════════════════════════════
% SECTION 8: CROSS-REFERENCES
% ═══════════════════════════════════════════════════════════

\subsubsection{Cross-References}

\begin{itemize}
\item Geometric derivation: Book 1, Chapter [X], Section [Y]
\item Related results: Sections~\ref{sec:[related1]}, \ref{sec:[related2]}
\item Open problems: OPR-[number] (Appendix~\ref{app:opr})
\item Baseline data: Table~\ref{tab:baseline_constants}
\end{itemize}

% ============================================================
% END OF TEMPLATE
% ============================================================

% ============================================================
% NOTES FOR AUTHORS:
% ============================================================
% 1. This template is MANDATORY for all major results
% 2. Delete instruction comments (% lines starting with "Example:")
% 3. Fill ALL [PLACEHOLDERS] with actual content
% 4. Verify all cross-references (\ref{...}) are correct
% 5. Ensure error budget numbers are realistic
% 6. Use appropriate epistemic tags [Dc:Sym/Num/Approx]
% 7. Status box should use tcolorbox package
% ============================================================
